\documentclass[11pt]{article}

%==================== Packages ====================
\usepackage[utf8]{inputenc}
\usepackage[T1]{fontenc}
\usepackage[margin=1in]{geometry}
\usepackage{amsmath,amssymb,amsthm}
\usepackage{mathtools}
\usepackage{algorithm}
\usepackage{algpseudocode}
\usepackage{enumitem}
\usepackage{hyperref}
\hypersetup{colorlinks=true,linkcolor=blue,citecolor=blue,urlcolor=blue}

%==================== Theorem environments ====================
\newtheorem{theorem}{Theorem}
\newtheorem{lemma}{Lemma}
\newtheorem{corollary}{Corollary}
\newtheorem{proposition}{Proposition}
\theoremstyle{definition}
\newtheorem{definition}{Definition}
\newtheorem{assumption}{Assumption}
\theoremstyle{remark}
\newtheorem{remark}{Remark}

%==================== Macros ====================
\newcommand{\OPT}{\mathrm{OPT}}
\newcommand{\OPTfrac}{\OPT_{\mathrm{frac}}}
\newcommand{\OPTSF}{\OPT_{\mathrm{SF}}}
\newcommand{\ALG}{\mathrm{ALG}}
\newcommand{\ALGSP}{\ALG_{\mathrm{SP}}}
\newcommand{\ALGSF}{\ALG_{\mathrm{SF}}}
\newcommand{\ESF}{E_{\mathrm{SF}}}
\newcommand{\src}{\mathrm{src}}
\newcommand{\dst}{\mathrm{dst}}
\newcommand{\dem}{\mathrm{dem}}
\newcommand{\len}{\mathrm{len}}
\newcommand{\Tin}{T^{(i)}_{\mathrm{in}}}
\newcommand{\Tout}{T^{(i)}_{\mathrm{out}}}
\newcommand{\din}{d_{\mathrm{in}}}
\newcommand{\dout}{d_{\mathrm{out}}}
\newcommand{\Din}{D^{(i)}_{\mathrm{in}}}
\newcommand{\Dout}{D^{(i)}_{\mathrm{out}}}
\DeclareMathOperator{\SP}{SP}

\title{\bf A Simplified MAL-C Problem:\\[2pt]
A $\theta$-Threshold Dual-Track Flow-Splitting Routing Algorithm and Its Analysis}
\author{}
\date{}

\begin{document}
\maketitle

\begin{abstract}
We study a simplified variant of the \emph{Minimum Aged Labelling with Capacity constraints}
(MAL-C) problem. The objective is to assign departure-time labels to the edges of a directed
graph and to route a set of freight demands so as to minimize the total number of unit-capacity
trucks dispatched, where the flow on each (edge, time) pair must be rounded \emph{up} to a whole
truck. We present a dual-track algorithm that routes bulky integer flow along shortest paths and
consolidates fractional flow onto a Directed Steiner Forest (DSF) backbone whose departure times
are produced by a global \emph{dual-label} construction. Under a store-and-forward model (freight
may wait at nodes at no cost), we prove that the backbone schedule is temporally feasible and that
every backbone edge carries at most two distinct departure labels, so that fractional flows from
different demands genuinely share trucks. Balancing the two tracks through a splitting threshold
$\theta$ yields an approximation ratio of $O\!\left(\sqrt{h}\,N^{1/3} + N^{2/3}\right)$, where
$N=|V|$ and $h$ is the number of demands.
\end{abstract}

%====================================================================
\section{Problem Description}
%====================================================================

We consider a simplified version of the \emph{Minimum Aged Labelling with Capacity constraints}
(MAL-C) problem. To isolate the core mechanisms---topological path selection and the construction
of time labels---we drop the original problem's \emph{maximum-label constraint} and its
\emph{demand-release times}.

\paragraph{Input.}
\begin{itemize}[topsep=2pt,itemsep=1pt]
  \item A directed graph $G=(V,E)$ together with an edge travel-time function
        $w:E\to\mathbb{N}$. We write $N=|V|$.
  \item A set of $h$ demands $D=\{d_1,\dots,d_h\}$. Each demand $d_i$ specifies a source
        $\src(d_i)$, a destination $\dst(d_i)$, and a freight amount $\dem(d_i)\in\mathbb{R}_{\ge 0}$.
\end{itemize}

\paragraph{Output and objective.}
The algorithm must, for every edge $e\in E$, decide a set of \emph{departure times} (the label set
$\lambda(e)$), and it must route every unit of freight from its source to its destination. Trucks
have capacity normalized to $1$. Even when the total flow that traverses an edge at a given time is
fractional, an entire truck must be dispatched (\emph{unconditional round-up}). Writing
$f_i(e,t)$ for the flow of demand $i$ on edge $e$ departing at time $t$, the goal is, subject to
flow conservation and connectivity, to minimize the total number of trucks:
\begin{equation}\label{eq:objective}
  \min \;\sum_{e\in E}\;\sum_{t}\;
  \left\lceil \sum_{i=1}^{h} f_i(e,t) \right\rceil .
\end{equation}

%====================================================================
\section{Algorithm Design}
%====================================================================

In a logistics network the largest source of waste is ``dispatching a truck that is not full.''
Our strategy is a \emph{dual-track} scheme. Bulky integer flow is routed directly along shortest
paths (SP) to save travel time, while sparse fractional flow---which would otherwise force every
demand to dispatch its own near-empty truck---is consolidated onto a backbone network, a Directed
Steiner Forest (DSF), where many demands ride together.

A splitting threshold $\theta\in(0,1)$ governs how much fractional flow is diverted to the
backbone. The main routine (Algorithm~\ref{alg:dualtrack}) calls the label-construction subroutine
(Algorithm~\ref{alg:duallabels}).

\begin{algorithm}[t]
\caption{$\theta$-Threshold Dual-Track Flow Splitting}
\label{alg:dualtrack}
\begin{algorithmic}[1]
\Require Graph $G=(V,E)$, demand set $D$ with $|D|=h$, splitting threshold $\theta\in(0,1)$.
\State Collect the sources and destinations of all demands; compute an approximate Directed Steiner
       Forest on $G$, obtaining the backbone $\ESF\subseteq E$.
\State $\lambda_{\mathrm{SF}}\gets \textsc{Assign-Dual-Labels}(\ESF,w)$
       \Comment{assign departure times to the backbone}
\For{each demand $d_i\in D$}
  \State Split $\dem(d_i)$ into an integer part $d_I$ and a fractional part $d_f$ with $0\le d_f<1$.
  \State \textbf{Integer routing:} send the $d_I$ flow along a shortest path of $G$, with the
         corresponding labels.
  \State \textbf{Fractional routing:}
  \If{$d_f \le 2\theta$}
     \State route the entire fractional flow $d_f$ on $\ESF$ (consolidated routing).
  \Else \Comment{$d_f > 2\theta$}
     \State force a flow of size $\theta$ onto $\ESF$ (consolidated routing);
     \State route the remaining $d_f-\theta$ on a shortest path of $G$.
  \EndIf
\EndFor
\State \textbf{Forced re-routing rule:} every fractional flow placed on $\ESF$ must strictly follow
       the \emph{time-varying simple path} produced by \textsc{Assign-Dual-Labels}: first travel up
       the reverse in-tree to a meeting point, then transfer to the forward out-tree.
\end{algorithmic}
\end{algorithm}

To ensure that the fractional flows on the backbone do not produce additional wasteful round-ups
merely because their departure times are misaligned, we design the label-construction subroutine so
that, by construction, \emph{every backbone edge carries at most two distinct departure labels}.

\begin{algorithm}[t]
\caption{Global Dual-Label Construction}
\label{alg:duallabels}
\begin{algorithmic}[1]
\Function{Assign-Dual-Labels}{$\ESF,w$}
  \State Contract $\ESF$ into a DAG and obtain its strongly connected components (SCCs)
         $C_1,C_2,\dots,C_K$ in topological order.
  \For{each SCC $C_i$}
     \State Choose a center $h_i$. Using the travel times $w(e)$, build, rooted at $h_i$:
     \State \quad a reverse \emph{in-tree} $\Tin$, recording the weighted depth $\din(\cdot)$ from a
            node to the root;
     \State \quad a forward \emph{out-tree} $\Tout$, recording the weighted depth $\dout(\cdot)$ from
            the root to a node.
     \State Let $\Din=\max_x \din(x)$ and $\Dout=\max_x \dout(x)$ be the in- and out-radii of $C_i$.
  \EndFor
  \State Compute a global base time $T_i$ for each SCC by dynamic programming, so that $T_i$ exceeds
         the schedule of every earlier SCC and dominates the arrival time of every cross-boundary
         edge entering $C_i$ (see Lemma~\ref{lem:cross}).
  \State \textbf{Label assignment.} Initialize $\lambda(e)=\emptyset$ for all $e\in\ESF$. For each
         edge $e=(x,y)\in\ESF$:
  \State \quad if $e\in\Tin$:\quad $t_{\mathrm{in}}(e)=T_i+\Din-\din(x)$;
  \State \quad if $e\in\Tout$:\quad $t_{\mathrm{out}}(e)=T_i+\Din+\dout(x)$;
  \State \quad if $e$ is a cross-boundary edge:\quad $t_{\mathrm{cross}}(e)=T_i+\Din+\Dout$.
  \State \Return the label set $\lambda_{\mathrm{SF}}$.
\EndFunction
\end{algorithmic}
\end{algorithm}

%====================================================================
\section{Temporal Feasibility and the Two-Label Guarantee}\label{sec:feasibility}
%====================================================================

The benefit of the dual-label construction relies on a single modeling assumption, which is
consistent with the objective~\eqref{eq:objective} because that objective charges only for trucks
\emph{on edges} and never for freight resting at a node.

\begin{assumption}[Free store-and-forward]\label{ass:wait}
Freight may be buffered at any node for an arbitrary amount of time at no cost. (This is admissible
precisely because the simplified problem omits the maximum-label / aging constraint of full MAL-C.)
\end{assumption}

The point of the construction is that each edge's label is a function of its \emph{position}
(its depth in a tree and the base time of its SCC), and is \emph{independent of which demand}
traverses it. We now show this yields a globally consistent schedule, so that all demands sharing an
edge depart on the same labels and therefore share trucks.

\begin{lemma}[Intra-SCC synchronization]\label{lem:intra}
Fix an SCC $C_i$. Consider any flow that travels up the in-tree $\Tin$ and then down the out-tree
$\Tout$ rooted at the center $h_i$. Under Assumption~\ref{ass:wait}, on every in-tree edge $e=(x,y)$
the flow arrives at $x$ no later than $t_{\mathrm{in}}(e)$, and on every out-tree edge it arrives no
later than $t_{\mathrm{out}}(e)$. Moreover the arrival time at any node depends only on the node, not
on the source of the flow.
\end{lemma}

\begin{proof}
\emph{In-tree.} Let $x_0\to x_1\to\cdots\to h_i$ be the root-ward path of a flow, with edges
$e_k=(x_k,x_{k+1})$. By definition of weighted depth, $\din(x_k)=w(e_k)+\din(x_{k+1})$. The flow
leaves $x_k$ at $t_{\mathrm{in}}(e_k)=T_i+\Din-\din(x_k)$, hence reaches $x_{k+1}$ at
\[
  \big(T_i+\Din-\din(x_k)\big)+w(e_k)
  = T_i+\Din-\big(\din(x_k)-w(e_k)\big)
  = T_i+\Din-\din(x_{k+1})
  = t_{\mathrm{in}}(e_{k+1}).
\]
Thus the arrival time at $x_{k+1}$ equals the departure label of its outgoing in-tree edge: the flow
proceeds with no waiting and, in particular, arrives at $x_{k+1}$ at the time
$T_i+\Din-\din(x_{k+1})$, which depends only on $x_{k+1}$. A flow that enters the tree at a deeper
node simply waits at its source until that node's outgoing label fires (Assumption~\ref{ass:wait}).
At the root the arrival time is $T_i+\Din-\din(h_i)=T_i+\Din$, independent of the source.

\emph{Out-tree.} Let $h_i\to y_1\to\cdots$ be a root-out path with edges $g_k=(y_k,y_{k+1})$ and
$\dout(y_{k+1})=\dout(y_k)+w(g_k)$. The first out-edge departs $h_i$ at
$t_{\mathrm{out}}(g_0)=T_i+\Din+\dout(h_i)=T_i+\Din$, exactly the arrival time at the root. Inductively
the flow reaches $y_{k+1}$ at $T_i+\Din+\dout(y_{k+1})=t_{\mathrm{out}}(g_{k+1})$, again synchronized
and source-independent.
\end{proof}

\begin{lemma}[Cross-boundary domination]\label{lem:cross}
There exist finite base times $T_1<T_2<\cdots<T_K$ such that for every cross-boundary edge
$e=(x,y)$ with $x\in C_i$ and $y\in C_j$ (so $i$ precedes $j$ in topological order), the flow reaches
$x$ no later than $t_{\mathrm{cross}}(e)=T_i+\Din+\Dout$, and reaches $y$ no later than the start of
$C_j$'s schedule. Consequently the whole route from any source to any destination is temporally
feasible under Assumption~\ref{ass:wait}.
\end{lemma}

\begin{proof}
By Lemma~\ref{lem:intra} a flow that has descended the out-tree of $C_i$ to the boundary node $x$
arrives at time $T_i+\Din+\dout(x)\le T_i+\Din+\Dout=t_{\mathrm{cross}}(e)$; it then waits until
$t_{\mathrm{cross}}(e)$ and crosses. It reaches $y$ at $t_{\mathrm{cross}}(e)+w(e)$. Process the SCCs
in topological order and set
\[
  T_j \;\ge\; \max_{e=(x,y):\,y\in C_j}\Big(t_{\mathrm{cross}}(e)+w(e)\Big),
\]
which is well defined because the contracted graph is a DAG and every such $e$ originates in an
earlier SCC. With this choice every flow entering $C_j$ has already arrived before $C_j$'s in-tree
schedule begins, so it waits at the entry node and joins the local schedule. Feasibility follows by
induction on the topological order.
\end{proof}

\begin{theorem}[Two labels per edge and consolidation]\label{thm:twolabels}
\textsc{Assign-Dual-Labels} assigns at most two distinct departure labels to every backbone edge, and
under Assumption~\ref{ass:wait} all fractional flows routed on a common backbone edge depart on those
labels. Hence each backbone edge incurs at most two round-ups in the objective~\eqref{eq:objective}.
\end{theorem}

\begin{proof}
Within an SCC the in-tree and out-tree are each a spanning arborescence of distinct orientation; an
edge $e=(x,y)$ may lie on the root-ward in-tree path \emph{and} on a root-out path, receiving the two
labels $t_{\mathrm{in}}(e)$ and $t_{\mathrm{out}}(e)$, but no more. A cross-boundary edge is disjoint
from the trees and receives the single label $t_{\mathrm{cross}}(e)$. Hence at most two labels per
edge. By Lemmas~\ref{lem:intra}--\ref{lem:cross} every flow using an edge arrives no later than the
relevant label and (with free waiting) departs exactly on it, so all demands sharing the edge are
consolidated onto its $\le 2$ labelled trucks.
\end{proof}

%====================================================================
\section{Approximation-Ratio Analysis}\label{sec:analysis}
%====================================================================

We compare the algorithm's cost $\ALG=\ALGSP+\ALGSF$ against the true optimum $\OPT$ of
\eqref{eq:objective}. Although $\OPT$ is unknown, it is bounded below by two ideal quantities.

\begin{definition}[Lower bounds]\label{def:lb}
\leavevmode
\begin{itemize}[topsep=2pt,itemsep=2pt]
  \item \textbf{Fractional-flow bound $\OPTfrac$:} the cost when fractional trucks are permitted (no
        round-up) and all flow travels along shortest paths,
        $\OPTfrac=\sum_{i=1}^{h}\dem(d_i)\cdot\len(\SP_i)$,
        where $\len(\SP_i)$ is the number of edges on a shortest $\src(d_i)$--$\dst(d_i)$ path.
        Since every unit of flow is forced to round up in reality, $\OPT\ge\OPTfrac$.
  \item \textbf{Connectivity bound $\OPTSF$:} regardless of how small a flow is, each source must be
        physically connected to its destination, so the optimum solution contains a Directed Steiner
        Forest; with unit edge cost, the minimum DSF edge count is a floor, $\OPT\ge\OPTSF$.
\end{itemize}
\end{definition}

Throughout, $\rho_{\mathrm{DSF}}=O(N^{2/3})$ denotes the approximation factor of the DSF routine used
in line~1 of Algorithm~\ref{alg:dualtrack}, so the backbone satisfies
$|\ESF|\le \rho_{\mathrm{DSF}}\cdot\OPTSF=O(N^{2/3})\cdot\OPTSF$.

%--------------------------------------------------------------------
\subsection{Cost of the shortest-path (SP) track}
%--------------------------------------------------------------------

The SP track carries (i) the integer parts $d_I$ of every demand and (ii) the leftover fractional
pieces $x=d_f-\theta$ for demands with $d_f>2\theta$.

\paragraph{Integer flow.} Integer flow never wastes a truck: on any edge the rounded value of a sum
of integers equals the sum itself, $\lceil\sum d_I\rceil=\sum d_I$. Hence
\begin{equation}\label{eq:spint}
  \ALGSP^{\mathrm{int}}=\sum_{i=1}^{h} d_I^{(i)}\cdot\len(\SP_i)
  \;\le\;\sum_{i=1}^{h}\dem(d_i)\cdot\len(\SP_i)=\OPTfrac .
\end{equation}

\paragraph{Leftover fractional flow.} By the splitting rule, a fractional piece reaches the SP track
only when $d_f>2\theta$, in which case its size satisfies the key inequality
\begin{equation}\label{eq:xtheta}
  x=d_f-\theta \;>\; 2\theta-\theta=\theta .
\end{equation}
On the SP track each such piece is rounded up to $\lceil x\rceil=1$ truck per edge. Using
$x>\theta\Rightarrow 1<x/\theta$ from \eqref{eq:xtheta},
\begin{equation}\label{eq:spfrac}
  \ALGSP^{\mathrm{frac}}
  =\sum_{i:\,d_f>2\theta} 1\cdot\len(\SP_i)
  \;<\;\sum_{i:\,d_f>2\theta}\frac{x_i}{\theta}\,\len(\SP_i)
  =\frac1\theta\sum_{i:\,d_f>2\theta} x_i\,\len(\SP_i)
  \;\le\;\frac1\theta\,\OPTfrac ,
\end{equation}
where the last step uses $x_i\le\dem(d_i)$ and that the pieces $\{x_i\}$ form a sub-collection of the
total flow, so $\sum x_i\len(\SP_i)\le\sum_i\dem(d_i)\len(\SP_i)=\OPTfrac$. (This is the corrected
form of the bound: the relation to $\OPTfrac$ is an \emph{inequality}, not an equality, and the
integer-flow contribution~\eqref{eq:spint} is accounted for separately.)

Adding \eqref{eq:spint} and \eqref{eq:spfrac},
\begin{equation}\label{eq:sp}
  \boxed{\;\ALGSP \;\le\; \Big(1+\tfrac1\theta\Big)\OPTfrac
  \;=\; O\!\Big(\tfrac1\theta\Big)\cdot\OPT.\;}
\end{equation}
Intuitively: the smaller $\theta$ is, the lower the threshold for diverting flow to the SP track, and
the larger the round-up blow-up factor $1/\theta$.

%--------------------------------------------------------------------
\subsection{Cost of the Steiner-forest (SF) track}
%--------------------------------------------------------------------

By the splitting rule, every demand injects at most $2\theta$ of flow into the backbone (it injects
$d_f\le 2\theta$ when $d_f\le 2\theta$, and exactly $\theta$ when $d_f>2\theta$). Let demand $i$ route
its backbone flow $s_i\le 2\theta$ along a path $R_i\subseteq\ESF$, and let $F_e=\sum_{i:\,e\in R_i}s_i$
be the total flow on backbone edge $e$.

By Theorem~\ref{thm:twolabels} each backbone edge carries at most two labels with flows
$f_{t_1}(e),f_{t_2}(e)$ summing to $F_e$. Since $\lceil z\rceil\le z+1$ for $z\ge 0$, the truck count
on a single edge obeys
\begin{equation}\label{eq:peredge}
  \lceil f_{t_1}(e)\rceil+\lceil f_{t_2}(e)\rceil \;\le\; F_e+2 .
\end{equation}
Summing \eqref{eq:peredge} over all backbone edges,
\begin{equation}\label{eq:sfsum}
  \ALGSF \;\le\; \sum_{e\in\ESF}\big(F_e+2\big)
  \;=\;\underbrace{\sum_{e\in\ESF}F_e}_{\text{flow term}}
  \;+\;\underbrace{2\,|\ESF|}_{\text{infrastructure term}} .
\end{equation}

\paragraph{Flow term.} Each backbone path has at most $|\ESF|$ edges, so
\begin{equation}\label{eq:flowterm}
  \sum_{e\in\ESF}F_e=\sum_{i=1}^{h} s_i\,|R_i|
  \;\le\; 2\theta\sum_{i=1}^{h}|R_i|
  \;\le\; 2\theta\cdot h\cdot|\ESF|
  \;=\; 2h\theta\,|\ESF| .
\end{equation}

\paragraph{Infrastructure term.} Even an edge that carries only an infinitesimal flow must dispatch a
whole truck per label; with at most two labels this is at most $2$ trucks per edge, giving the
$2|\ESF|$ term. \emph{This term is independent of $\theta$ and must not be dropped:} it is the
unavoidable cost of physically operating the backbone.

Combining \eqref{eq:sfsum}--\eqref{eq:flowterm} and $|\ESF|\le O(N^{2/3})\OPTSF$,
\begin{equation}\label{eq:sf}
  \boxed{\;\ALGSF \;\le\; 2(h\theta+1)\,|\ESF|
  \;=\; O\!\big((h\theta+1)\,N^{2/3}\big)\cdot\OPTSF
  \;\le\; O\!\big((h\theta+1)\,N^{2/3}\big)\cdot\OPT.\;}
\end{equation}
Intuitively: the larger $\theta$ is, the more total flow you allow into the backbone and the more the
flow term grows; the additive $N^{2/3}$ is the price of building the backbone at all.

%--------------------------------------------------------------------
\subsection{Balancing the bounds and the final ratio}
%--------------------------------------------------------------------

Adding \eqref{eq:sp} and \eqref{eq:sf}, the overall ratio is
\begin{equation}\label{eq:total}
  \mathrm{TotalRatio}(\theta)
  \;=\; O\!\Big(\tfrac1\theta\Big)
  \;+\; O\!\big(h\theta\,N^{2/3}\big)
  \;+\; O\!\big(N^{2/3}\big).
\end{equation}
Only the first two terms depend on $\theta$. Minimizing $\tfrac1\theta+h\theta N^{2/3}$ by equating
the two,
\begin{equation}\label{eq:thetastar}
  \frac1\theta = h\theta\,N^{2/3}
  \;\Longrightarrow\; \theta^2=\frac{1}{h\,N^{2/3}}
  \;\Longrightarrow\; \theta^\ast=\frac{1}{\sqrt{h}\,N^{1/3}}\in(0,1).
\end{equation}
At $\theta^\ast$ each of the first two terms equals $\sqrt{h}\,N^{1/3}$, while the third is unchanged,
so
\begin{equation}\label{eq:final}
  \boxed{\;\mathrm{TotalRatio}(\theta^\ast)
  \;=\; O\!\big(\sqrt{h}\,N^{1/3}+N^{2/3}\big)
  \;=\; O\!\big(\max\{\sqrt{h}\,N^{1/3},\,N^{2/3}\}\big).\;}
\end{equation}

\begin{theorem}[Approximation guarantee]\label{thm:main}
Under Assumption~\ref{ass:wait}, Algorithm~\ref{alg:dualtrack} with threshold
$\theta^\ast=1/(\sqrt{h}\,N^{1/3})$ produces a feasible routing and labelling whose cost is within a
factor $O\!\big(\sqrt{h}\,N^{1/3}+N^{2/3}\big)$ of the optimum of~\eqref{eq:objective}.
\end{theorem}

\begin{remark}[Two regimes]
The bound \eqref{eq:final} separates into two regimes, with the crossover at $h=N^{2/3}$:
\[
  \mathrm{TotalRatio}(\theta^\ast)=
  \begin{cases}
    O\!\big(\sqrt{h}\,N^{1/3}\big), & h\ge N^{2/3}\ \ (\text{round-up dominated}),\\[4pt]
    O\!\big(N^{2/3}\big),           & h< N^{2/3}\ \ (\text{backbone dominated}).
  \end{cases}
\]
When demands are plentiful the dual-track splitting is what limits the cost; when demands are few the
cost is governed by the cost of the Directed Steiner Forest itself, i.e.\ by the DSF approximation
factor $\rho_{\mathrm{DSF}}=O(N^{2/3})$. No choice of $\theta$ removes the $N^{2/3}$ term, since the
backbone must be built to enable consolidation.
\end{remark}

\begin{remark}[On a common pitfall]
A naive analysis drops the infrastructure term $2|\ESF|$ in \eqref{eq:sfsum}, folding it into
$O(h\theta\,N^{2/3})$. That is valid only when $h\theta\ge 1$. At the balancing optimum, however,
$h\theta^\ast=\sqrt{h}/N^{1/3}<1$ whenever $h<N^{2/3}$, so the additive term in fact \emph{dominates}
and cannot be absorbed. Retaining it yields the honest bound \eqref{eq:final}; omitting it would
incorrectly claim $O(\sqrt{h}\,N^{1/3})$.
\end{remark}

%====================================================================
\section{Conclusion}
%====================================================================

We presented a dual-track routing-and-labelling algorithm for the simplified MAL-C problem. The
$\theta$-threshold rule sends bulky integer flow along shortest paths and consolidates fractional
flow onto a Directed Steiner Forest backbone. The global dual-label construction equips the backbone
with a temporally consistent schedule (Theorem~\ref{thm:twolabels}): under a free store-and-forward
model, intra-component routing is exactly synchronized and cross-component routing is made feasible
by base times that dominate all incoming arrivals, so that each backbone edge carries at most two
labels and fractional flows from different demands genuinely share trucks. A careful, complete cost
analysis---retaining the $\theta$-independent backbone term---gives an approximation ratio of
\[
  O\!\big(\sqrt{h}\,N^{1/3}+N^{2/3}\big),
\]
attained at $\theta^\ast=1/(\sqrt{h}\,N^{1/3})$. The bound is round-up dominated when $h\ge N^{2/3}$
and backbone dominated otherwise, the latter reflecting the intrinsic cost of the Directed Steiner
Forest approximation.

\end{document}
